% Copyright (c) 2026 Marcio Luiz Lima de Oliveira and V. el Hakim.
% Distribution and modification are subject to LPPL version 1.3c.
% SPDX-License-Identifier: LPPL-1.3c
% Part of the LPPL-covered work listed in MANIFEST.txt.
% See LICENSE and LICENSING.txt for the licence and scope.

\chapter{Extend and maintain the layout}\label{chap:07}
\renewcommand{\FigLocation}{Chapters/Chapter_7/fig/}
\renewcommand{\ChapLocation}{Chapters/Chapter_7/}
\begin{absbox}
Keep structural changes local and predictable. This chapter explains document divisions, partial builds, the package responsibilities and the custom commands that connect the layout.
\end{absbox}

\section{Chapters and section files}
An ordinary chapter begins with this pattern:
\begin{lstlisting}
\chapter{A Descriptive Title}\label{chap:example}
\renewcommand{\FigLocation}{Chapters/Chapter_3/fig/}
\renewcommand{\ChapLocation}{Chapters/Chapter_3/}
\begin{absbox}
  A short account of this chapter's purpose and contribution.
\end{absbox}
\section{The first section}
\end{lstlisting}
\verb|\FigLocation| is a convenience prefix for graphics;
\verb|\ChapLocation| is a prefix for section files. They are ordinary macros,
not a change of working directory. Relative paths still resolve from the root
where \texttt{main.tex} is compiled. Reassign them in each chapter rather than
relying on whichever chapter happened to run before it.

\verb|\include| gives each chapter its own auxiliary file and starts a new page.
\verb|\input| inserts a smaller file without those chapter-level page boundaries.
Keep \verb|\include| calls in \texttt{main.tex} and use \verb|\input| for sections.
The original four section files in Chapter 2 demonstrate that distinction.

\section{Partial builds and back matter}
Use \verb|\includeonly{Chapters/Chapter_3/Chapter_03}| in the preamble for a
drafting shortcut. Build the complete thesis first so omitted chapters have
auxiliary records, and build it completely again before checking lists and
citations. The custom colour counter and symbol-use state are not a guarantee
of final output in a partial build; do not judge first-use pages from one.

\verb|\frontmatter| uses front-matter numbering;
\verb|\mainmatter| starts numbered chapters;
\verb|\backmatter| suppresses chapter numbering again. The list of symbols is
printed by \verb|\PrintThesisSymbols| inside \texttt{MathSymbols.tex}. Keep the
label \texttt{ch:listofsymbols} there: symbol hyperlinks target that label.
\verb|\ThesisSymbolRecord| and \verb|\ThesisUntrackedAcronyms| are supporting
registry commands, not commands to insert into normal prose. The former reads
auxiliary records; the latter prevents glossary descriptions from changing
acronym first-use state.

To add a numbered appendix, put its source inside the existing
\texttt{Chapter\_Appendix} directory and add it before \verb|\backmatter|:
\begin{lstlisting}
\appendix
\include{Chapters/Chapter_Appendix/AdditionalMethods}
\backmatter
\end{lstlisting}
Its file begins with \verb|\chapter{Additional methods}|. Numbered appendices
also exercise the chapter style, so extend the marker count/palette or disable
markers for that part. The existing abbreviation and notation files are
unnumbered back-matter chapters and do not require \verb|\appendix|.

\section{Package responsibilities}
Load packages once in the preamble. The groups below help locate a setting
without treating every package's internal commands as part of this template.
\begin{description}
\item[Page and text design] \texttt{memoir}, \texttt{geometry}, \texttt{titlesec},
\texttt{microtype}, \texttt{fontenc}, \texttt{lmodern}, \texttt{setspace} and
\texttt{footmisc} provide the book structure, dimensions, fonts and spacing.
\item[Colour and graphics] \texttt{xcolor}, \texttt{graphicx}, \texttt{tikz},
\texttt{tcolorbox}, \texttt{subfig}, \texttt{wrapfig} and \texttt{pdfpages}
support artwork, boxes and imported pages. \texttt{circuitikz} is available for
circuit diagrams; its \texttt{siunitx} option supplies the unit tools used here.
The inherited \texttt{background} package is neutralised by an empty background;
the example covers explicitly draw their own page overlays.
\item[Mathematics] \texttt{amsmath}, \texttt{amsthm}, \texttt{amssymb} and
\texttt{mathtools} support displays, proofs, symbols and paired delimiters.
\item[Tables and lists] \texttt{booktabs}, \texttt{tabularx}, \texttt{longtable},
\texttt{multirow}, \texttt{colortbl}, \texttt{multicol} and \texttt{enumitem}
provide rules, wrapping columns, multipage tables and list controls.
\texttt{varwidth} and \texttt{minibox} support compact text boxes.
\item[Language and references] \texttt{babel}, \texttt{csquotes},
\texttt{acronym}, \texttt{biblatex} and \texttt{hyperref} handle language,
quotations, abbreviation state, citations and links. Only Dutch and English
are enabled in this template. \texttt{alphabeta} supports Greek letters in text.
\item[Code and algorithms] \texttt{listings} prints literal source examples;
\texttt{algorithm}, \texttt{algorithmicx} and \texttt{algpseudocode} provide the
pseudocode example. \texttt{ifthen} selects optional cover artwork;
\texttt{intcalc} supports integer arithmetic in inherited design code.
\end{description}
The inherited \texttt{lipsum} and \texttt{breakcites} dependencies are not needed
for these examples and have been omitted. Consult the class and package manuals
before adding another package that manages the same captions, chapter styles
or bibliography. The template does not need a shell-escape-enabled build.

\section{Keep links and first-use state understandable}
\verb|\hypersetup| controls PDF metadata and link colours. The configuration
sets acronym links black through public font hooks; the symbol commands also
use black links. Setting \texttt{linkcolor} therefore changes ordinary
cross-references without recolouring every tracked symbol or abbreviation.
The wrappers preserve the acronym package's fonts through three command copies:
\begin{lstlisting}
\ThesisOriginalAcsFont
\ThesisOriginalAcfFont
\ThesisOriginalAcfsFont
\end{lstlisting}
Leave these copies intact when only changing a palette.

Use meaningful text in bookmarks, avoid fragile formatting in captions, and
give figures a short list entry when the full caption is long. Cross-references
are records built across runs, not values calculated correctly in every first
pass. After changing a bibliography backend or moving many labels, rebuild
the document's auxiliary files through your editor's clean/recompile operation.
